% ============================================================
%  SAVI: Conditional Viterbi Integration
%  Revised Technical Design Memo — Companion to the SAVI Draft
% ============================================================

\documentclass[11pt,letterpaper]{article}
\usepackage{times}
\usepackage[margin=1.25in]{geometry}
\usepackage{amsmath,amssymb,amsthm}
\usepackage{booktabs}
\usepackage{algorithm}
\usepackage{algpseudocode}
\usepackage[hidelinks]{hyperref}
\usepackage{microtype}
\usepackage{tabularx}
\usepackage{caption}
\usepackage{enumitem}
\usepackage{mdframed}
\usepackage{fancyhdr}
\usepackage{xspace}

\pagestyle{fancy}
\fancyhf{}
\fancyhead[L]{\small\itshape SAVI: Conditional Viterbi Integration}
\fancyhead[R]{\small\itshape Revised Technical Design Memo}
\fancyfoot[C]{\thepage}
\renewcommand{\headrulewidth}{0.4pt}
\setlength{\headheight}{14pt}

% ── Theorem environments ─────────────────────────────────────────────────────
\newtheorem{definition}{Definition}
\newtheorem{proposition}{Proposition}
\newtheorem{theorem}{Theorem}
\newtheorem{corollary}{Corollary}
\newtheorem{remark}{Remark}

% ── Math macros ──────────────────────────────────────────────────────────────
\newcommand{\Phifn}{\Phi}
\newcommand{\Pfwd}{P_{\mathrm{fwd}}}
\newcommand{\Pbwd}{P_{\mathrm{bwd}}}
\newcommand{\calS}{\mathcal{S}}
\newcommand{\calP}{\mathcal{P}}
\newcommand{\SAVI}{\textsc{savi}\xspace}
\newcommand{\eps}{\varepsilon}

% ── Callout box ──────────────────────────────────────────────────────────────
\newmdenv[
  linecolor=black,
  linewidth=0.8pt,
  innerleftmargin=10pt,
  innerrightmargin=10pt,
  innertopmargin=8pt,
  innerbottommargin=8pt,
  skipabove=8pt,
  skipbelow=8pt
]{callout}

% ─────────────────────────────────────────────────────────────────────────────
\title{\textbf{Integrating Viterbi Dynamic Programming into \SAVI}\\
       \large Canonical Prompting, Backpointers, and Conditional Trellis Inference\\[6pt]
       \normalsize Revised Technical Design Memo --- Companion to the SAVI Draft}
\author{}
\date{}

% =============================================================================
\begin{document}
\maketitle
\thispagestyle{fancy}

\begin{abstract}
The current \SAVI draft correctly frames semantic-state search as Viterbi-inspired, because ordinary
LLM prompting conditions on the full text history and therefore violates the Markov property required
by classical Viterbi inference.
This memo integrates the useful parts of the Viterbi upgrade while avoiding an overclaim.
The central change is \emph{canonical prompting}: after merging histories that reach the same semantic
state, the next LLM query is reconstructed from a deterministic representation of that state rather than
from an arbitrary preserved text history.
Under this induced process, \SAVI can be implemented as a conditional Viterbi dynamic program over
a beam-limited semantic-state trellis, with explicit backpointers and path recovery.
The result is not unqualified exact Viterbi over the original LLM text process.
It is exact only under stated conditions: sound semantic mapping, complete candidate support or
accepted candidate truncation, no beam-pruning error, sufficient state expressiveness, and a calibrated
learned value function for the bidirectional extension.
The revised positioning is therefore stronger and safer:
\SAVI is Viterbi-inspired in its general form;
\emph{Canonical} \SAVI implements conditional Viterbi DP over the induced semantic-state trellis.
\end{abstract}

\begin{callout}
\textbf{Executive summary.}
Do not replace every occurrence of ``Viterbi-inspired'' with ``exact Viterbi.''
Instead, introduce \textbf{Canonical \SAVI} as the principled variant.
Canonical prompting makes the engineered inference process Markovian;
DP backpointers make path recovery clean;
and the correctness theorem should be conditional, not absolute.
\end{callout}

\tableofcontents
\newpage

% ── 1  What Should Change ────────────────────────────────────────────────────
\section{What Should Change in the SAVI Paper}

\subsection{The Safe Claim}

The upgraded claim should be:

\begin{quote}
\itshape
\SAVI is Viterbi-inspired in its general form.
Under canonical prompting, it becomes a conditional Viterbi dynamic program over an induced
semantic-state trellis, exact up to candidate truncation, beam pruning, semantic-map soundness,
state-representation expressiveness, and learned value approximation.
\end{quote}

This is strong enough to improve the theory and safe enough to survive review.

\subsection{Three Variants to Define}

The paper should explicitly distinguish three variants.

\vspace{4pt}
\begin{center}
\small
\begin{tabularx}{\linewidth}{lXX}
\toprule
\textbf{Variant} & \textbf{Prompt source} & \textbf{Claim} \\
\midrule
History \SAVI
  & Preserved representative text history.
  & Practical version. Viterbi-inspired, but Markov approximation may be nonzero. \\
\addlinespace
Canonical \SAVI
  & Deterministic prompt built from the canonical semantic state.
  & Conditional Viterbi DP over the induced semantic-state process. \\
\addlinespace
Bidirectional Canonical \SAVI
  & Canonical state prompt plus learned goal-conditioned value.
  & Forward Viterbi DP plus a learned approximation to the backward variable. \\
\bottomrule
\end{tabularx}
\end{center}
\vspace{4pt}

This resolves the main risk in the earlier memo: it preserves the Viterbi upgrade without claiming that
the original history-conditioned LLM process is exactly Markovian.

% ── 2  Why Base SAVI Is Viterbi-Inspired ────────────────────────────────────
\section{Why the Original \SAVI Is Only Viterbi-Inspired}

Classical Viterbi inference requires a state abstraction that is sufficient for future transitions.
More concretely, it requires:

\begin{enumerate}[label=(\roman*)]
  \item \textbf{A state space.} There must be states $S_t \in \calS$.
  \item \textbf{A transition process.} Future states must be induced by transitions from the current
        state.
  \item \textbf{A Markov property.} The distribution over future transitions must depend on the
        current state, not on the full history.
\end{enumerate}

The original \SAVI draft already has the first two ingredients: the semantic map $\Phifn$ constructs
states and candidate chunks induce transitions.
The critical gap is the Markov property.
LLMs condition on text histories.
If two text histories $H_A$ and $H_B$ satisfy $\Phifn(H_A) = \Phifn(H_B)$, the LLM may still
generate differently from them because the visible prompts differ.

\begin{definition}[History-conditioned Markov error]
For merged histories $H_A, H_B$ with $\Phifn(H_A) = \Phifn(H_B)$, define
\[
  \eps(H_A, H_B) = D\bigl[\Pfwd(\cdot \mid H_A) \;\|\; \Pfwd(\cdot \mid H_B)\bigr],
\]
where $D$ may be KL, JS divergence, or an empirical divergence over induced next semantic states.
\end{definition}

When $\eps$ is large, merging is lossy: the algorithm discards one history even though that history
may have produced better future continuations.
This is why the base method should remain described as Viterbi-inspired.

% ── 3  Canonical Prompting ───────────────────────────────────────────────────
\section{Canonical Prompting: The Key Integration}

\begin{definition}[Canonical prompting]
A prompting strategy is \emph{canonical} if the LLM input at step $t$ is a deterministic function of
the current semantic state and the fixed problem statement:
\[
  \mathrm{prompt}_t = f(S_{t-1}, \calP),
\]
independent of the text history used to reach $S_{t-1}$.
\end{definition}

Under canonical prompting, two histories that merge into the same state produce the same next prompt.
Thus, for the induced inference process,
\[
  \Pfwd(x_t \mid H_A) = \Pfwd(x_t \mid H_B)
  \quad \text{whenever} \quad \Phifn(H_A) = \Phifn(H_B),
\]
provided both histories are replaced by the same canonical prompt.
The Markov violation is not solved for the original text process; rather, the algorithm deliberately
constructs a new state-conditioned process.

\subsection{Canonical Prompt Templates}

\begin{itemize}
  \item \textbf{Symbolic algebra.}
        The prompt contains the original problem, the current canonical SymPy expression or equation,
        declared assumptions, and a fixed instruction to propose the next legal algebraic transformation.
        Prior derivation text is omitted.
  \item \textbf{Restricted code generation.}
        The prompt contains the function signature, problem specification, bounded observational state,
        visible output traces on fixed public inputs, and a fixed instruction to propose the next code
        chunk.
        Prior arbitrary prose is omitted or summarized into the canonical execution state.
  \item \textbf{Theorem proving.}
        The prompt contains the formal proof state, local hypotheses, goals, and allowed tactics.
        Prior natural-language proof text is not used as state.
\end{itemize}

\begin{callout}
\textbf{Tradeoff.}
Canonical prompting removes path dependence, but it can lose useful context if the semantic state
representation is too poor.
This is not a theoretical footnote; it must be tested.
The paper should add an experiment comparing canonical-state prompting against preserved-history
prompting.
\end{callout}

% ── 4  Viterbi DP with Backpointers ─────────────────────────────────────────
\section{Viterbi DP with Backpointers}

\subsection{Forward Recurrence}

Let $V_t(S)$ be the best accumulated log score of any path reaching canonical state $S$ at step $t$.
With canonical prompting, define
\[
  V_0(S_0) = 0, \qquad V_0(S) = -\infty \text{ for } S \neq S_0.
\]
The recurrence is
\begin{equation}
  V_t(S') = \max_{\substack{S \in \mathcal{B}_{t-1} \\ x:\,\Phifn(S,x) = S'}}
    \Bigl[V_{t-1}(S) + \log \Pfwd\bigl(x \mid f(S, \calP)\bigr)\Bigr].
  \label{eq:viterbi-dp}
\end{equation}

The backpointer is
\[
  \pi_t(S') = \arg\max_{(S,x):\,\Phifn(S,x) = S'}
    \Bigl[V_{t-1}(S) + \log \Pfwd\bigl(x \mid f(S, \calP)\bigr)\Bigr].
\]

\subsection{Bidirectional Extension}

For Bidirectional Canonical \SAVI, add a learned state value:
\begin{equation}
  V_t^\lambda(S') = \max_{S,x:\,\Phifn(S,x) = S'}
    \Bigl[V_{t-1}^\lambda(S) + \log \Pfwd(x \mid f(S,\calP))
      + \lambda \log \Pbwd(G{=}1 \mid S', \calP)\Bigr].
  \label{eq:bidir-dp}
\end{equation}

This is not the exact HMM backward variable unless $\Pbwd$ is perfectly calibrated to the induced
transition process.
The safe statement is that $\Pbwd$ is a learned approximation to a backward value.

\begin{definition}[Backward-value approximation]
$\Pbwd(G{=}1 \mid S_t, \calP)$ estimates the probability that state $S_t$ can reach a correct
terminal solution under the induced decoding policy and task objective.
It is a \emph{learned value function}, not a future-token predictor.
\end{definition}

\subsection{Algorithm}

\begin{algorithm}[H]
\caption{Canonical \SAVI as Conditional Viterbi DP}
\label{alg:canonical-savi}
\begin{algorithmic}[1]
  \Require Problem $\calP$, semantic map $\Phifn$, LLM $\Pfwd$, optional value $\Pbwd$,
           candidates $N$, beam width $K$, steps $T$, weight $\lambda$
  \State Initialize $V_0(S_0) \leftarrow 0$, $\mathcal{B}_0 \leftarrow \{S_0\}$
  \For{$t = 1, \ldots, T$}
    \State $\mathcal{C} \leftarrow \emptyset$
    \For{each $S \in \mathcal{B}_{t-1}$}
      \State $p \leftarrow f(S, \calP)$  \Comment{canonical prompt}
      \State Generate $N$ candidate chunks $x_1, \ldots, x_N \sim \Pfwd(\cdot \mid p)$
      \For{each candidate $x_i$}
        \State $S' \leftarrow \Phifn(S, x_i)$
        \If{$S' \neq \bot$}
          \State $s \leftarrow V_{t-1}(S) + \log \Pfwd(x_i \mid p)
                   + \lambda \log \Pbwd(G{=}1 \mid S', \calP)$
          \State Add $(S',\, s,\, S,\, x_i)$ to $\mathcal{C}$
        \EndIf
      \EndFor
    \EndFor
    \For{each unique state $S'$ in $\mathcal{C}$}
      \State $V_t(S') \leftarrow \max\{s : (S', s, S, x) \in \mathcal{C}\}$
      \State $\pi_t(S') \leftarrow \arg\max\{s : (S', s, S, x) \in \mathcal{C}\}$
    \EndFor
    \State $\mathcal{B}_t \leftarrow$ top-$K$ states by $V_t(\cdot)$
  \EndFor
  \State $S^* \leftarrow \arg\max_{S \in \mathcal{B}_T} V_T(S)$
  \State \Return trajectory recovered by backtracking through $\pi_T, \ldots, \pi_1$
\end{algorithmic}
\end{algorithm}

The key implementation changes are modest: prompt from canonical state, store backpointers, and
recover the final trajectory by backtracking rather than by returning the last retained history.

% ── 5  Conditional Correctness Theorem ──────────────────────────────────────
\section{Conditional Correctness Theorem}

\begin{theorem}[Conditional Viterbi correctness]
\label{thm:conditional-correctness}
Assume the following conditions for the induced canonical-prompting process:
\begin{enumerate}[label=(\roman*)]
  \item \textbf{Canonical prompting:} all prompts are deterministic functions $f(S_t, \calP)$.
  \item \textbf{Sound semantic map:} $\Phifn$ has no false merges over the evaluated domain.
  \item \textbf{Complete candidate support:} the candidate generator enumerates every transition with
        nonzero probability, or else the theorem is restricted to the sampled candidate set.
  \item \textbf{No beam-pruning error:} the optimal path's state remains within the top-$K$ retained
        states at every step.
  \item \textbf{(Bidirectional \SAVI only)} $\Pbwd$ is a calibrated estimate of terminal success
        probability under the induced process.
\end{enumerate}
Then Algorithm~\ref{alg:canonical-savi} with $\lambda = 0$ recovers the maximum-score path in the
induced semantic-state trellis.
With $\lambda > 0$, it recovers the maximum-score path under the augmented forward-plus-value
objective in Equation~\eqref{eq:bidir-dp}.
\end{theorem}

\begin{proof}[Proof sketch]
Under canonical prompting, future candidate distributions depend only on the current canonical state.
Under semantic-map soundness, merging preserves equivalence classes rather than collapsing distinct
states.
Equation~\eqref{eq:viterbi-dp} is therefore the standard Viterbi max recurrence over the induced
trellis.
Backpointers recover the maximizing path.
Candidate truncation and beam pruning restrict the trellis on which the theorem applies; if the
optimal path is pruned or never sampled, no beam method can recover it.
The bidirectional case is exact only for the stated augmented objective and depends on the calibration
quality of $\Pbwd$.
\end{proof}

\begin{remark}
This theorem is intentionally conditional.
That is a feature, not a weakness.
Each condition corresponds to an empirical audit: false-merge rate, candidate budget, beam-width
sensitivity, canonical-prompting validation, and value calibration.
\end{remark}

% ── 6  How to Modify the Paper Text ─────────────────────────────────────────
\section{How to Modify the Paper Text}

\subsection{Abstract and Contributions}

Replace any unqualified claim of ``exact Viterbi'' with this language:

\begin{quote}
\itshape
\SAVI is Viterbi-inspired in its general form.
We further introduce Canonical \SAVI, which reconstructs the LLM prompt from the canonical
semantic state after each merge.
This induces a state-conditioned process on which \SAVI can be implemented as a conditional Viterbi
dynamic program with explicit backpointers and path recovery.
\end{quote}

Add a contribution bullet:

\begin{quote}
\itshape
We define Canonical \SAVI and prove a conditional correctness theorem showing that, under sound
semantic mapping, sufficient candidate support, and no beam-pruning error, the algorithm recovers the
maximum-score path in the induced semantic-state trellis.
\end{quote}

\subsection{Framework Section}

Add a new subsection after the semantic map definitions:

\begin{quote}
\textbf{Canonical prompting.}
The base form of \SAVI may continue from the representative text history associated with each retained
state.
This is practical but only approximately Markovian.
Canonical \SAVI instead constructs the next prompt from the canonical state itself.
This enforces the beam invariant not only at the state level but also at the prompting level:
all paths that merge into the same state produce the same future prompt.
\end{quote}

Then replace the recurrence with Equation~\eqref{eq:viterbi-dp} and introduce the backpointer
table $\pi_t$.

\subsection{Theory Section}

Do not delete the Markov approximation discussion.
Split it:

\begin{enumerate}
  \item \textbf{History \SAVI:} Markov approximation error must be measured because prompts differ
        across merged histories.
  \item \textbf{Canonical \SAVI:} Markov error is removed by construction for the induced process,
        but the burden shifts to state expressiveness and semantic-map soundness.
\end{enumerate}

This is better than claiming the Markov problem vanishes universally.

\subsection{Bidirectional \SAVI Section}

Replace ``$\Pbwd$ is the true backward variable'' with:

\begin{quote}
\itshape
$\Pbwd$ is a learned approximation to the backward value: the probability that a canonical state can
reach a correct terminal solution under the induced decoding process.
If calibrated, it plays the role of the HMM backward variable; otherwise it should be treated as a
value heuristic.
\end{quote}

% ── 7  Experiments to Add ───────────────────────────────────────────────────
\section{Experiments to Add}

The Viterbi integration needs experiments that specifically test whether the theory helps, not just
whether the original semantic merging idea works.

\vspace{4pt}
\begin{center}
\small
\begin{tabularx}{\linewidth}{lXX}
\toprule
\textbf{Experiment} & \textbf{Comparison} & \textbf{Purpose} \\
\midrule
E6a Canonical prompt validation
  & History \SAVI vs.\ Canonical \SAVI.
  & Tests whether state-only prompting loses useful information.
    Success: canonical accuracy within 2--3 percentage points, or better after
    enriching the state representation. \\
\addlinespace
E6b Markov reduction
  & Next-state distributions from two merged histories before and after canonical prompting.
  & Shows that canonical prompting collapses history-dependent divergence in the induced process. \\
\addlinespace
E6c DP backtracking ablation
  & Beam-with-merge vs.\ Viterbi DP with backpointers.
  & Tests whether principled path recovery improves performance, especially on longer chains. \\
\addlinespace
E6d State expressiveness ablation
  & Minimal state prompt vs.\ enriched state prompt.
  & Identifies whether canonical prompting fails because the state representation omits
    necessary context. \\
\addlinespace
E6e Value calibration
  & Reliability curve and expected calibration error for $\Pbwd$.
  & Determines whether Bidirectional \SAVI can legitimately be interpreted as using a
    backward-value approximation. \\
\bottomrule
\end{tabularx}
\end{center}
\vspace{4pt}

These should be added to the experimental section after the existing aliasing and
online-vs-post-hoc experiments.
They do not replace BAR, CMR, RSD, false-merge audits, or iso-compute comparisons.

% ── 8  Revised Nomenclature ─────────────────────────────────────────────────
\section{Revised Nomenclature}

\begin{center}
\small
\begin{tabularx}{\linewidth}{lX}
\toprule
\textbf{Avoid} & \textbf{Use instead} \\
\midrule
``\SAVI is exact Viterbi''
  & ``Canonical \SAVI is conditional Viterbi DP over the induced semantic-state trellis.'' \\
\addlinespace
``Remove Viterbi-inspired everywhere''
  & ``Base \SAVI is Viterbi-inspired; Canonical \SAVI is the Viterbi-DP variant.'' \\
\addlinespace
``$\Pbwd$ is the true backward variable''
  & ``$\Pbwd$ is a learned approximation to the backward value, calibrated against terminal
    success.'' \\
\addlinespace
``Markov error collapses to zero''
  & ``Canonical prompting removes history dependence in the induced process; residual risk is
    state expressiveness and $\Phifn$ soundness.'' \\
\addlinespace
``Viterbi optimality''
  & ``Optimality within the sampled, beam-limited, soundly merged semantic trellis.'' \\
\bottomrule
\end{tabularx}
\end{center}

% ── 9  Risk Analysis ────────────────────────────────────────────────────────
\section{Risk Analysis}

\begin{center}
\small
\begin{tabularx}{\linewidth}{lXX}
\toprule
\textbf{Risk} & \textbf{Problem} & \textbf{Mitigation} \\
\midrule
Canonical prompting loses context
  & The semantic state may omit information contained in the original text history.
  & Run E6a and E6d. Enrich the state prompt with assumptions, variable bindings, constraints,
    and selected verified provenance if needed. \\
\addlinespace
False merges corrupt DP
  & Viterbi DP assumes that merged paths are genuinely equivalent.
  & Keep existing false-merge audit; report false-merge and missed-merge rates separately
    for algebra and code. \\
\addlinespace
Candidate truncation weakens exactness
  & Sampling only $N$ chunks per state means the full transition distribution is not enumerated.
  & State that the theorem applies to the sampled trellis; sweep $N$ to measure sensitivity. \\
\addlinespace
Beam pruning drops optimal states
  & Top-$K$ pruning can eliminate the globally best path.
  & Sweep $K$ and report performance/coverage curves. \\
\addlinespace
$\Pbwd$ is uncalibrated
  & An uncalibrated value model can harm Bidirectional \SAVI.
  & Add E6e; report calibration and perform $\lambda$ sensitivity.
    Keep Forward \SAVI as the main contribution. \\
\bottomrule
\end{tabularx}
\end{center}

% ── 10  Recommended Integration Order ───────────────────────────────────────
\section{Recommended Integration Order}

\begin{enumerate}
  \item Add Canonical \SAVI as a named variant, not as a replacement for all \SAVI.
  \item Modify the algorithm to include canonical prompting, DP table updates, and backpointers.
  \item Add the conditional correctness theorem with explicit assumptions.
  \item Add the experiments E6a--E6e to validate canonical prompting and DP backtracking.
  \item Soften the $\Pbwd$ language to ``learned approximation to the backward value.''
  \item Keep Forward \SAVI as the main empirical contribution until $\Pbwd$ is validated.
\end{enumerate}

\begin{callout}
\textbf{Final recommended sentence for the paper.}
\SAVI converts text-space tree search into semantic-state trellis search.
In its practical history-conditioned form it is Viterbi-inspired;
in its canonical-prompting form it implements conditional Viterbi DP over the induced
semantic-state trellis, with approximation sources made explicit and empirically measurable.
\end{callout}

\end{document}
